
I enjoy
teaching mathematics, from first-year calculus to more advanced courses, and see teaching and dissemination as a central part of mathematics.
In teaching mathematics,
I aim to give students not only a clear and rigorous foundation in the {material}, but also the intuitive understanding that is so hard to learn from
books alone.
{I try to build this intuition through visualizations and concrete examples, and to always connect it back to precise statements and proofs.}

During my time at Cambridge,
I have {supervised} for the courses Algebraic Topology,
Riemann Surfaces, and Complex Analysis.
Supervisions are Cambridge's version of an exercise class.
In lieu of going through the problems in a standard classroom format,
students meet a supervisor in small groups of two or three.
While the format might be peculiar,
it has taught me a great deal.
Through such direct contact with the students,
I was able to develop a good sense of which parts of the courses
are easy to follow, and which ones introduce new ideas that
require the student to go through a paradigm shift in their own thinking.
Such steps take much longer for students to understand.
For instance, in Riemann Surfaces, almost every student struggled with the step from treating $\log z$ as a ``multivalued function'' to seeing it as an honest function on a surface. Drawing the surface sheet by sheet and walking around the branch point together helped them over this hurdle.
Similarly, in Algebraic Topology, students found it hard to connect the abstract definition of homology to the intuitive picture of ``holes in the space''; going through some explicit computations with simplicial chain complexes cleared this up.
While each student is different and learns in their own way and at their own pace,
I was able to see patterns across all students in my supervisions.
When teaching any undergraduate course on related topics in the future,
I will be able to draw on the deeper insight I gained while supervising at Cambridge.

Even before starting my PhD, I was a teaching assistant during my Master's and
undergraduate studies and gained experience in traditional classroom
teaching.
At ETH Zurich, I taught exercise
classes for a series of
Analysis courses,
ranging from Analysis II to Analysis IV: a sequence running from multivariable calculus, through measure theory and Lebesgue integration, to Hilbert spaces, Fourier analysis, and the spectral theorem.
Analysis II and III are mandatory courses, while Analysis IV is aimed at students with a deeper interest in analysis.
Teaching exercise classes for this sequence of courses of increasing difficulty and
specialization was a rewarding experience.
I especially enjoyed seeing some students progress all the way through these classes,
and witnessing their mathematical growth.

In my exercise classes, I made sure to focus on the conceptual understanding of the students.
To ensure that everyone was always up to speed on the material from the lectures, I encouraged students
to submit questions about the previous lectures ahead of the class
and did a review session
at the beginning of each class.
This also gave students who did not want to ask in front of the whole class a way to raise their difficulties.
In addition to analyzing the submitted questions to identify any stumbling blocks,
I used my own experience of learning this material during my undergraduate degree.
My aim was always to give the students a geometric and intuitive understanding of the material.
For this purpose, I prepared computer visualizations for key concepts such as the Weierstrass function and the Fourier modes.
These visual elements were very well received by the students, and I plan to keep using them in my teaching.
To make the classes interactive and create an environment in which students were happy
to ask questions and speak up,
I usually ended each class by discussing an easy unseen problem and solving it together with the students.
This also allowed me to gauge how well the class had understood the material, and helped me plan the next session and give the lecturer feedback on the students' progress.


I am looking forward to teaching my own courses.
During my PhD, I gained my first experience of designing a course,
when I taught the preparatory class in algebraic topology for incoming Master's students at Cambridge in 2024.
Such preparatory classes, consisting of two sessions and a problem sheet,
are meant to ensure that all incoming students are well prepared for the Master's degree.
Having worked closely with the lecturers at ETH, I have also seen how a course is run from the inside,
and I plan to bring the methods described above into my own lectures:
collecting questions ahead of time, using visualizations for key concepts,
and regularly working through short problems together with the students.
As a former teaching assistant, I know how important it is for the lecturer and the teaching assistants to stay in close contact,
and I will make sure that the teaching assistants in my courses get the feedback and support they need.
I would be happy to teach calculus, linear algebra, and real analysis, as well as more advanced courses in topology, complex analysis, and differential geometry.
At the graduate level, I would like to offer a topics course in symplectic geometry, leading up to the geometry of Hamiltonian diffeomorphism groups.
I would also like to supervise undergraduate research and reading projects.
Parts of my research have accessible computational aspects, which could form the basis of good projects for undergraduates.

Whether in a supervision with two students or in a lecture course, my aim remains the same: to give students both a rigorous foundation and the intuition that makes the material their own.
